\chapter{Minimale, maximale, kleinste und größte Elemente}
\label{b12guide:ch:extremal}

\section{Begriffe und Charakterisierungen}

\begin{definition}[Extremalelemente einer Teilmenge]
\label{b12guide:def:extremal}
Für $T\subseteq P$ seien
\[
\begin{aligned}
 \MinP(T)&=\{m\in T\mid\forall t\in T\;(t\leq m\Rightarrow t=m)\},\\
 \MaxP(T)&=\{m\in T\mid\forall t\in T\;(m\leq t\Rightarrow t=m)\}.
\end{aligned}
\]
Ihre Elemente heißen \emph{minimal} beziehungsweise \emph{maximal}
in $T$. Ein $m\in T$ heißt \emph{Minimum} oder \emph{kleinstes
Element}, wenn $m\leq t$ für alle $t\in T$ gilt.
Ein $M\in T$ heißt \emph{Maximum} oder \emph{größtes Element},
wenn $t\leq M$ für alle $t\in T$ gilt.
Falls sie existieren, werden sie mit $\min T$ und $\max T$ bezeichnet.
Die Symbole $\MinP(T)$ und $\MaxP(T)$ bezeichnen dagegen stets Mengen.
Aussagen wie $m=\min T$ schließen die Existenz des bezeichneten
Elements ein.
\end{definition}

\begin{theorem}[Strenge Kriterien für Minimalität und Maximalität]
Für $m\in T\subseteq P$ gilt
\[
\begin{aligned}
 m\in\MinP(T)&\quad\Longleftrightarrow\quad
                    \neg\exists t\in T\;(t<m),\\
 m\in\MaxP(T)&\quad\Longleftrightarrow\quad
                    \neg\exists t\in T\;(m<t).
\end{aligned}
\]
Für $T=\varnothing$ sind beide Mengen leer; Minimum und Maximum
existieren dann nicht.
\end{theorem}

\begin{theorem}[Eindeutigkeit und Verhältnis der Extremalbegriffe]
\label{b12guide:thm:least}
Jede Teilmenge $T\subseteq P$ besitzt höchstens ein Minimum und
höchstens ein Maximum. Es gilt
\[
 \min T=m\ \Rightarrow\ \MinP(T)=\{m\},\qquad
 \max T=M\ \Rightarrow\ \MaxP(T)=\{M\}.
\]
Die Existenz genau eines minimalen beziehungsweise maximalen
Elements genügt im allgemeinen nicht für die Existenz eines
Minimums beziehungsweise Maximums. Auch bei $T\ne\varnothing$
müssen minimale oder maximale Elemente nicht existieren.
\end{theorem}

\begin{theorem}[Vergleichbarkeit eines extremalen Elements]
\label{b12guide:thm:extremal-comparable}
Für $m\in\MinP(T)$ gilt
\[
 m=\min T\quad\Longleftrightarrow\quad
       \forall t\in T\;(m\leq t\ \lor\ t\leq m).
\]
Für $M\in\MaxP(T)$ gilt entsprechend
\[
 M=\max T\quad\Longleftrightarrow\quad
       \forall t\in T\;(M\leq t\ \lor\ t\leq M).
\]
Dabei wird nur die Vergleichbarkeit mit dem genannten Element
verlangt, keine Totalität der Ordnung auf $T$.
\end{theorem}

\section{Antiketten und Vererbung}

\begin{definition}[Antikette]\label{b12guide:def:antichain}
Eine Teilmenge $C\subseteq P$ heißt \emph{Antikette}, wenn
je zwei verschiedene Elemente von $C$ unvergleichbar sind.
Die leere Menge und jede einelementige Teilmenge sind Antiketten.
\end{definition}

\begin{theorem}[Extremalmengen sind Antiketten]
\label{b12guide:thm:extremal-antichains}
Für jede Teilmenge $T\subseteq P$ sind $\MinP(T)$ und $\MaxP(T)$
Antiketten. Ferner sind äquivalent:
\[
 T\text{ ist eine Antikette},\qquad
 \MinP(T)=T,\qquad \MaxP(T)=T.
\]
\end{theorem}

\begin{theorem}[Vererbung auf Teilmengen]
\label{b12guide:thm:extremal-subsets}
Für $S\subseteq T\subseteq P$ gilt
\[
 \MinP(T)\cap S\subseteq\MinP(S),\qquad
 \MaxP(T)\cap S\subseteq\MaxP(S).
\]
Ist $m=\min T$ und $m\in S$, so ist $m=\min S$.
Ist $M=\max T$ und $M\in S$, so ist $M=\max S$.
Die umgekehrten Inklusionen gelten im allgemeinen nicht.
\end{theorem}

\begin{theorem}[Zusätzliche Vergleiche und Extremalelemente]
\label{b12guide:thm:extremal-refinement}
Seien $R,S\in\PO(P)$ mit $R\subseteq S$. Für jedes $T\subseteq P$
gilt, mit einem Index zur Angabe der jeweiligen Ordnung,
\[
 \MinP_S(T)\subseteq\MinP_R(T),\qquad
 \MaxP_S(T)\subseteq\MaxP_R(T).
\]
Ein Minimum oder Maximum bezüglich $R$ bleibt ein Minimum
beziehungsweise Maximum bezüglich $S$.
\end{theorem}

\begin{theorem}[Dualität der Extremalbegriffe]
Für jedes $T\subseteq P$ gilt
\[
 \MinP_{P^{\mathrm{op}}}(T)=\MaxP_P(T),\qquad
 \MaxP_{P^{\mathrm{op}}}(T)=\MinP_P(T).
\]
Minimum und Maximum werden ebenfalls vertauscht, einschließlich
ihrer jeweiligen Existenzbedingungen.
\end{theorem}
